\documentclass[10pt,a4paper]{amsart}
\usepackage[margin=19mm]{geometry}
\usepackage{amsmath,amssymb,mathtools}
\usepackage[scr=rsfs,cal=euler]{mathalfa}
\usepackage{xfrac}
\usepackage[unicode,hidelinks,bookmarks=false]{hyperref}
\newcommand{\ie}{i.e.\ }
\newcommand{\cf}{cf.\ }
\newcommand{\resp}{respectively\ }
\newcommand{\Subsection}{\subsection}
\providecommand{\nobreakdash}{\nobreak\hbox{-}}
\setlength{\emergencystretch}{2em}
\multlinegap0pt
\usepackage{bm}
\usepackage{bbm}
\usepackage{dsfont}
\usepackage{xspace}
\usepackage{cancel}
\usepackage{mathtools}
\usepackage{amsthm}
\makeatletter
\@ifundefined{theorem}{\newtheorem{theorem}{Theorem}}{}
\makeatother
\title[The Rigidity of Constraint: A Spencer-Hodge Theoretic Approa]{The Rigidity of Constraint: A Spencer-Hodge Theoretic Approach to the Hodge Conjecture}
\author{Dongzhe Zheng}
\date{Source: arXiv:2506.12720v1 (CC BY 4.0)}
\begin{document}
\maketitle

\noindent\emph{Source and licence.} The Rigidity of Constraint: A Spencer-Hodge Theoretic Approach to the Hodge Conjecture. Version arXiv:2506.12720v1 (CC BY 4.0), distributed under the Creative Commons Attribution 4.0 International licence, as recorded in the arXiv licence metadata for this exact version and shown on the abstract page https://arxiv.org/abs/2506.12720v1. Full rights notice: licenses/arxiv-2506.12720-CC-BY-4.0.txt.\par\smallskip

\noindent\textit{Editorial scope.} Counted unit: the Theorem whose source label is \texttt{thm:nilpotency} in the article (source0.tex lines 2726--2729, source label \texttt{thm:nilpotency}), delivered together with the article's own complete original proof of it (source0.tex lines 2731--2826, one \texttt{proof} environment of 96 lines). No mathematics is written, repaired or re-derived: statement and proof are the article's own text line for line. Required context retained uncounted: none. Every definition, display and label of those lines is retained unchanged. Omitted from this package and not redistributed: 1--2725, 2730--2730, 2827--3280. Nothing inside a retained excerpt is omitted or altered: every retained body is delivered exactly as the article's lines are written, so no omission receipt of any kind accompanies this package. Citations: the retained excerpts carry no citation command at all, so no bibliography entry of the article is transcribed and no reference is redistributed. Reference adaptation: none was needed and none was made; every reference inside the delivered unit resolves inside it, so no unresolved reference or citation is produced. Printed slips kept literal: no printed word, symbol, hypothesis or number of the counted statement or of its proof was altered. Layout and class: the article's own class is \texttt{extarticle} with packages including inputenc, CJKutf8, geometry, graphicx, booktabs, mathrsfs, bm, xcolor, algorithm, algpseudocode, mathtools, enumitem, tikz-cd, newtxtext; the wrapper is the lane's pinned \texttt{amsart} editorial wrapper at 10pt on A4, which restates the theorem-like environments the retained text needs and adds the article's own macro definitions that the retained text needs (none), with extra wrapper packages (bm, bbm, dsfont, xspace, cancel, mathtools). The delivered numbering is the unit's own. Assets: no retained span contains a figure environment, a graphics-inclusion command, a PDF-page inclusion, a table or a TikZ picture; no image file is copied into this package, the package declares no asset dependency and no image file is redistributed with it. Licence: version arXiv:2506.12720v1 of this article is distributed under the Creative Commons Attribution 4.0 International licence, as recorded in the arXiv licence metadata for this exact version and shown on the abstract page https://arxiv.org/abs/2506.12720v1. A full rights notice is delivered at licenses/arxiv-2506.12720-CC-BY-4.0.txt. Characters: every retained excerpt is pure ASCII, so the delivered engine, XeLaTeX with its default Latin Modern text font, reports no missing character.
\begin{theorem}[Nilpotency]
\label{thm:nilpotency}
$(\delta^\lambda_\mathfrak{g})^2 = 0$.
\end{theorem}

\begin{proof}
We proceed by induction on tensor degree.

\textbf{Base case}: $k=1$. Need to prove that for any $v \in \mathfrak{g}$, $\delta^\lambda_\mathfrak{g}(\delta^\lambda_\mathfrak{g}(v)) = 0$.

Let $\sigma = \delta^\lambda_\mathfrak{g}(v) \in \operatorname{Sym}^2(\mathfrak{g})$. Let $\{H_1, \ldots, H_n\}$ be basis of $\mathfrak{g}$, then:
$$\sigma = \sum_{1 \leq i \leq j \leq n} \sigma_{ij} (H_i \odot H_j)$$
where $\sigma_{ij} = \sigma(H_i, H_j) = (\delta^\lambda_\mathfrak{g}(v))(H_i, H_j)$.

Applying Leibniz rule:
\begin{align}
\delta^\lambda_\mathfrak{g}(\sigma) &= \sum_{1 \leq i \leq j \leq n} \sigma_{ij} \delta^\lambda_\mathfrak{g}(H_i \odot H_j) \\
&= \sum_{1 \leq i \leq j \leq n} \sigma_{ij} [\delta^\lambda_\mathfrak{g}(H_i) \odot H_j + H_i \odot \delta^\lambda_\mathfrak{g}(H_j)]
\end{align}

For test vectors $(w_1, w_2, w_3)$:
\begin{align}
(\delta^\lambda_\mathfrak{g}(\sigma))(w_1, w_2, w_3) &= \sum_{1 \leq i \leq j \leq n} \sigma_{ij} [(\delta^\lambda_\mathfrak{g}(H_i) \odot H_j)(w_1, w_2, w_3) \\
&\quad + (H_i \odot \delta^\lambda_\mathfrak{g}(H_j))(w_1, w_2, w_3)]
\end{align}

Now compute $(\delta^\lambda_\mathfrak{g}(H_i) \odot H_j)(w_1, w_2, w_3)$. Since this is 3rd order symmetric tensor, need to average over all permutations:
\begin{align}
&(\delta^\lambda_\mathfrak{g}(H_i) \odot H_j)(w_1, w_2, w_3) \\
&= \frac{1}{3!} \sum_{\pi \in S_3} (\delta^\lambda_\mathfrak{g}(H_i) \odot H_j)(w_{\pi(1)}, w_{\pi(2)}, w_{\pi(3)}) \\
&= \frac{1}{6} \sum_{\pi \in S_3} (\delta^\lambda_\mathfrak{g}(H_i))(w_{\pi(1)}, w_{\pi(2)}) \cdot B(H_j, w_{\pi(3)})
\end{align}

where $B$ is Killing form, $B(H_j, w) = \langle H_j^*, w \rangle$.

Now key is computing $\sigma_{ij} = (\delta^\lambda_\mathfrak{g}(v))(H_i, H_j)$:
\begin{align}
\sigma_{ij} &= \frac{1}{2} \left( \langle \lambda, [H_i, [H_j, v]] \rangle + \langle \lambda, [H_j, [H_i, v]] \rangle \right)
\end{align}

Let $v = \sum_k v_k H_k$, then:
\begin{align}
[H_j, v] &= \sum_k v_k [H_j, H_k] = \sum_{k,l} v_k C_{jk}^l H_l \\
[H_i, [H_j, v]] &= \sum_{k,l} v_k C_{jk}^l [H_i, H_l] = \sum_{k,l,m} v_k C_{jk}^l C_{il}^m H_m
\end{align}

where $C_{jk}^l$ are structure constants of Lie algebra.

Therefore:
\begin{align}
\sigma_{ij} &= \frac{1}{2} \sum_{k,l,m} v_k \left( C_{jk}^l C_{il}^m + C_{ik}^l C_{jl}^m \right) \langle \lambda, H_m \rangle
\end{align}

Now we need to prove that when we substitute these into $(\delta^\lambda_\mathfrak{g}(\sigma))(w_1, w_2, w_3)$ and expand all terms, result is zero.

Key observation: Final expression will contain terms of form $\langle \lambda, [w_a, [w_b, [w_c, v]]] \rangle$ (after symmetrization).

\textbf{Core lemma}: For any $v \in \mathfrak{g}$ and $w_1, w_2, w_3 \in \mathfrak{g}$, define:
$$T(v; w_1, w_2, w_3) = \sum_{\text{symmetrization}} \langle \lambda, [w_i, [w_j, [w_k, v]]] \rangle$$

where sum runs over all permutations $(i,j,k)$ of $(w_1, w_2, w_3)$.

We prove $T(v; w_1, w_2, w_3) = 0$.

\textbf{Detailed computation}:
\begin{align}
T(v; w_1, w_2, w_3) &= \langle \lambda, [w_1, [w_2, [w_3, v]]] \rangle + \langle \lambda, [w_1, [w_3, [w_2, v]]] \rangle \\
&\quad + \langle \lambda, [w_2, [w_1, [w_3, v]]] \rangle + \langle \lambda, [w_2, [w_3, [w_1, v]]] \rangle \\
&\quad + \langle \lambda, [w_3, [w_1, [w_2, v]]] \rangle + \langle \lambda, [w_3, [w_2, [w_1, v]]] \rangle
\end{align}

Using Jacobi identity $[X, [Y, Z]] + [Y, [Z, X]] + [Z, [X, Y]] = 0$:

For $[w_2, [w_3, v]]$:
$$[w_2, [w_3, v]] + [w_3, [v, w_2]] + [v, [w_2, w_3]] = 0$$
Therefore:
$$[w_2, [w_3, v]] = -[w_3, [v, w_2]] - [v, [w_2, w_3]]$$

Similarly:
$$[w_3, [w_2, v]] = -[w_2, [v, w_3]] - [v, [w_3, w_2]] = -[w_2, [v, w_3]] + [v, [w_2, w_3]]$$

Now:
\begin{align}
&\langle \lambda, [w_1, [w_2, [w_3, v]]] \rangle + \langle \lambda, [w_1, [w_3, [w_2, v]]] \rangle \\
&= \langle \lambda, [w_1, (-[w_3, [v, w_2]] - [v, [w_2, w_3]])] \rangle \\
&\quad + \langle \lambda, [w_1, (-[w_2, [v, w_3]] + [v, [w_2, w_3]])] \rangle \\
&= -\langle \lambda, [w_1, [w_3, [v, w_2]]] \rangle - \langle \lambda, [w_1, [w_2, [v, w_3]]] \rangle
\end{align}

Continuing to apply Jacobi identity to remaining terms, and using $[v, w] = -[w, v]$, we can systematically prove all terms cancel each other, yielding $T(v; w_1, w_2, w_3) = 0$.

\textbf{Inductive step}: Assume nilpotency holds for all degrees $\leq k$. For $s = s_1 \odot s_2$ (total degree $k+1$), using Leibniz rule:
\begin{align}
(\delta^\lambda_\mathfrak{g})^2(s_1 \odot s_2) &= \delta^\lambda_\mathfrak{g}[\delta^\lambda_\mathfrak{g}(s_1) \odot s_2 + (-1)^{\deg(s_1)} s_1 \odot \delta^\lambda_\mathfrak{g}(s_2)] \\
&= \delta^\lambda_\mathfrak{g}(\delta^\lambda_\mathfrak{g}(s_1)) \odot s_2 + (-1)^{\deg(s_1)+1} \delta^\lambda_\mathfrak{g}(s_1) \odot \delta^\lambda_\mathfrak{g}(s_2) \\
&\quad + (-1)^{\deg(s_1)} \delta^\lambda_\mathfrak{g}(s_1) \odot \delta^\lambda_\mathfrak{g}(s_2) + (-1)^{\deg(s_1)} s_1 \odot \delta^\lambda_\mathfrak{g}(\delta^\lambda_\mathfrak{g}(s_2))
\end{align}

First and last terms are zero by inductive hypothesis. Middle two terms:
$$(-1)^{\deg(s_1)+1} \delta^\lambda_\mathfrak{g}(s_1) \odot \delta^\lambda_\mathfrak{g}(s_2) + (-1)^{\deg(s_1)} \delta^\lambda_\mathfrak{g}(s_1) \odot \delta^\lambda_\mathfrak{g}(s_2) = 0$$
\end{proof}

\end{document}
